
\subsection{\ac{CPU} und Speicher Architektur}
Computerarchitekturen bestehen aus mehreren Ebenen zur
Abstraktion der Hardware, von einzelnen Transistoren bis hin
zu einem Betriebssystem und Benutzerprogrammen.
Die Funktionsweise von \ac{CbCC} und
Transient Execution befindet sich auf der Ebene der
\ac{CPU}-Architektur, ihrer Mikroarchitektur sowie der
darunter liegenden Hardware.
Eine Architektur ist die strukturelle Sicht des Computers, die
dem Softwareentwickler präsentiert wird.
Anstatt diese Architektur direkt in Hardware zu bauen und
Entwicklern zur Verfügung zu stellen, wird bei modernen
Prozessoren eine zusätzliche Abstraktionsschicht in Form von
Mikrocode eingefügt. Diese ermöglicht die Implementierung
von Instruktionen in Software und damit auch das nachträgliche
patchen von Bugs.
(\cite[14,379]{sarangi_basic_2021})
Dementsprechend sind für jede Schicht nur bedingt die Vorgänge
der darunter liegenden Schicht sichtbar. Sowohl \ac{CbCC} als
auch \acp{TEA}
nutzen Vorgänge in ihnen nicht zugänglichen, unteren Schichten
aus, um in der Regel Informationen aus anderen Software
Bereichen zu lesen. Aufgrund der Hardwarenähe der Schichten,
die ausgenutzt werden ist das Erkennen entsprechender
Informationsflüsse sehr schwer.

\subsubsection{Computer Instruktionen und Ressourcen}
Eine grundsätzliche Rechnerarchitektur ist heute als die
„Von-Neumann-Architektur“ bekannt. Die zentrale Komponente
dieser Architektur ist die \ac{CPU}, welche für reine
arithmetische Berechnungen (\ac{ALU}) und für die
Programminterpretation zuständig ist
(Control Unit). Des Weiteren enthält diese Architektur einem
Speicher (engl. Memory) sowie Geräte für Eingabe und
Ausgabe. Dieser Speicher wird auch Arbeitsspeicher
beziehungsweise \ac{RAM} genannt.
Bei modernen Computern ist der \ac{RAM}, aufgrund seiner Größe im
Verhältnis zur Geschwindigkeit der Rechenoperationen
vergleichsweise langsam.
Aus diesem Grund verwendet die \ac{CPU} selbst interne Register mit
einer Kopie der aktuell benötigten Werte aus dem \ac{RAM}.
(\cite[35]{sarangi_basic_2021})

Ein einzelner Computer kann mehrere Prozessoren besitzen, die
Programme parallel ausführen können. Wenn sich mehrere
Prozessoren auf demselben Chip befinden, werden diese Kerne
genannt. Mittlerweile gehört der größte Teil der Computersysteme
zu den Multicore-Systemen mit eigenen Speicherzugriffen.
(\cite[569,578]{sarangi_basic_2021})

Die unterste Softwareebene, die für den Entwickler bereitgestellt
wird, ist die \ac{ISA} einer \ac{CPU}. Diese beschreibt ursprünglich
die Befehle, die von der \ac{CPU} in Hardware ausgeführt werden
können. Während arithmetische Operationen dieser Instruktionen
direkt auf den Registern ausgeführt werden, wird die
Speicherarchitektur auf primär die Instruktionen „load“ und
„store“ abstrahiert. Zwar existieren weitere speicherrelevante
Instruktionen, diese hängen jedoch von der jeweiligen \ac{ISA}
und teilweise auch von der konkreten Hardwareimplementierung ab.
Bei modernen Prozessoren ist die \ac{ISA} außerdem nicht direkt
in der Hardware implementiert, sondern wird durch sogenannte
\acp{mop} abstrahiert. Da diese für die vorliegende
Arbeit nur von untergeordneter Bedeutung sind, werden sie in
Anhang \ref{appendix:microinstructions} näher
beschrieben. Entsprechend arbeitet die \ac{ISA} mit dem
architektonischen, während die Mikroinstruktionen mit dem
mikroarchitektonischen Zustand der \ac{CPU} arbeitet.
(\cite[19, 373, 379]{sarangi_basic_2021})
Auch wenn verschiedene \acp{ISA} existieren und die besprochenen
Konzepte auf die meisten Implementierungen von \acp{CPU}
zutreffen, wird in dieser Arbeit \textendash\space sofern nicht
anders angegeben \textendash\space primär die x86-64 \ac{ISA}
betrachtet. Im Folgenden ist außerdem mit Speicher stets der
nicht persistente \ac{RAM} des Computers gemeint, der von den
\ac{CPU} während der Ausführung verwendet wird.

Da sich bei der Von-Neumann-Architektur Instruktionen und Daten
im selben Speicher befinden, wird der Speicher überwiegend als
ein großes Array mit aufsteigenden Adressen betrachtet.
Adressiert wird dabei jeweils ein Byte. Eine Folge von Instruktionen,
die von einem Computer ausgeführt werden soll,
wird als Programm bezeichnet und befindet sich genauso wie die
Daten im Speicher.
Um ein Programm auszuführen, müssen die Instruktionen zunächst
von der \ac{CPU} geladen (engl. fetch), anschließend von der Hardware
interpretiert und schließlich ausgeführt werden.
Die Adresse der aktuell
auszuführenden Instruktion wird als Programmzähler
(engl. \acf{PC}) bezeichnet.
(\cite[74,188]{mandl_grundkurs_2020})

\subsubsection{Speicher und Cache Architektur}
Ein zentrales Problem moderner Computer ist die Abhängigkeit von
Zugriffen auf den \ac{RAM}. Die Speichergröße der \ac{CPU}-internen
Register ist begrenzt und reicht für komplexere Programmabläufe
nicht aus.
Deshalb befinden sich sowohl Programme als auch deren Werte
im Speicher und müssen durch „load“ und „store“ Instruktionen mit
den Registern synchronisiert werden. Dementsprechend besteht ein
erheblicher Anteil jedes Programms aus Speicher Zugriffen.
Allerdings ist die Menge an Speicher auf demselben Chip wie
der \ac{CPU} aus Platzgründen deutlich begrenzt.

Grundsätzlich lassen sich zwei für moderne Rechnersysteme
relevante \ac{RAM}-Tech\-no\-lo\-gien unterscheiden: \ac{SRAM} und
\ac{DRAM}.
An sich ist \ac{SRAM} die deutlich überlegene Technologie mit
signifikant geringeren Zugriffszeiten. Jedoch ist \ac{SRAM}
besonders bei der Anordnung in größeren Mengen erheblich größer
und teurer. Auch wenn die modernsten Server Prozessoren heute
bis zu einem Gigabyte \ac{SRAM} (\cite{noauthor_server_nodate})
auf dem Chip erreichen, ist dies besonders im Verhältnis zu der
Anzahl der Kerne des Prozessors sehr gering.
Der Speicher auf dem \ac{CPU}-Chip ist allerdings nur ein Cache,
der einen Teil des gesamten Speichers des Computers hält und
die Zugriffszeit für möglichst viele Zugriffe minimieren soll.
Dieser Cache ist nötig, da seit dem Ende des 20. Jahrhunderts
die Speicherzugriffszeit die Rechenzeit der \ac{CPU} deutlich
übersteigt.
(\cite[37,38]{handy_cache_1993}, \cite[13]{hat_what_2007})

Der restliche Programmspeicher befindet sich außerhalb der \ac{CPU}
und muss über einen Bus angesteuert werden. Während ein idealer
Speicher eine Zugriffszeit von einem \ac{CPU}-Zyklus hätte und
damit praktisch instantan wäre, ist die reale Zugriffszeit
abhängig von der Speicherebene, die angefragt wird. Diese Zeit
liegt in der Regel zwischen ein bis zwei und wenigen 100 Zyklen
(\cite[723,724]{yarom_flushreload_2014}),
je nachdem welche Speicher Schicht angefragt wird.
(\cite[16]{hat_what_2007}, \cite[534]{sarangi_basic_2021})

Da der Cache nur ein Abbild eines Teils des gesamten Speichers
ist, muss
dieser so implementiert werden, dass sich bei Bedarf beliebige
Adressen in den Cache schreiben und von dort lesen lassen.
Gleichzeitig muss ein effizienter Mechanismus implementiert
werden, um auf physischer Ebene anhand der Speicheradresse den
Wert der Adresse zu ermitteln. Ist eine Adresse bereits im Cache
vorhanden und kann direkt zurückgegeben werden wird dies ein
Cache „Hit“, andernfalls ein „Miss“ genannt
(\cite[514]{sarangi_basic_2021}). Das Entfernen
einer Adresse beziehungsweise eines Blocks wird „Eviction“
beziehungsweise als Verb „evict“ genannt.

Empirisch zeigt sich, dass dicht
nebeneinander liegende Adressen häufig im Zusammenhang verwendet
werden. Diese Beobachtung wird als räumliche Lokalität
bezeichnet. Aus diesem Grund speichern Caches Daten nicht
byteweise, sondern in Blöcken. Anstatt des
Ladens einer einzelnen Adresse wird also ein fester
Speicherbereich \textendash\space typischerweise mit einer Größe
zwischen 32 und 128 Byte \textendash\space in
den Cache geladen. Der Cache ermittelt anhand der zu ladenden
Adresse den gesuchten Block und gibt den darin gesuchten
Wert zurück. Um von der Speicheradresse auf den physischen Cache
Eintrag schließen zu können, wird eine Mappingmethode verwendet
(vgl. \ref{appendix:basic-cache-arch}).
In der Praxis wird häufig ein \ac{SAC} verwendet. Dabei ist
jede Speicheradresse einer begrenzten Menge von physischen
Cache-Einträgen, einem sogenannten „Set“, zugeordnet. Die Anzahl
an physischen Einträgen je Set wird häufig, wie beispielsweise
in Anhang \ref{appendix:lscpu-cache}, als „Ways“ bezeichnet.
Die Einträge eines Sets werden über eine Setadresse angesteuert,
die von der physischen Adresse abgeleitet wird. Dies kann
durch das Assoziieren eines Teils der Speicheradresse
mit der Setadresse geschehen. Zum Beispiel könnten die
ersten $N$ Bit einer Adresse die Setadresse darstellen. Dabei
wird das Mapping meist so gewählt, dass Adressen, die sich im
selben Set befinden, nicht direkt nebeneinander liegen. Denn
Adressen im selben Set konkurrieren gewissermaßen um den Cache
und nach dem Prinzip der räumlichen Lokalität werden Adressen,
die dicht beieinander liegen, häufig zur gleichen Zeit genutzt.
(\cite[515,516]{sarangi_basic_2021}, \cite[16]{handy_cache_1993})

Eine Konsequenz aus dieser Architektur ist die
Existenz von sogenannten Eviction-Sets. Ein Eviction-Set ist eine
Menge an physischen Adressen, die alle dem gleichen Cache-Set
zugeordnet sind. Da diese Adressen bewusst nicht nah beieinander
gewählt werden, ist es zudem nicht unwahrscheinlich, dass
verschiedene Programme virtuelle Adressen besitzen, deren
physische Adressen demselben Eviction-Set angehören.
(\cite[40,41]{vila_theory_2019})

Aufgrund der Größe der \ac{SRAM} Zellen des Cache nimmt mit
zunehmender Cache Größe auch die Zugriffszeit zu. Aus diesem
Grund ist der Cache in Prozessoren in mehrere Ebenen
mit unterschiedlichen Größen geteilt. Die meisten Prozessoren
haben 3 verschiedene Cache „Layer“, die mit L1, L2, L3 bezeichnet
werden. Der letzte Cache vor dem \ac{RAM} wird \ac{LLC} genannt.
(\cite[388]{slomka_grundlagen_2023}, \cite[596]{sarangi_basic_2021})

Um die logische Implementierung dieser Cache Layer zu
vereinfachen, kann ein Cache „inclusive“ gegenüber einer höheren
Schicht sein. Das heißt, dass alle Blöcke, die dieser Cache
enthält, sich auch in dem darüberliegenden Cache befindet.
Inclusive \ac{CPU}-Speichersysteme sind aufgrund der
einfacheren Implementierung deutlich weiter verbreitet. Einige
neuere Prozessoren nutzen mittlerweile allerdings non-inclusive
Cache Layer (\cite[513]{sarangi_basic_2021},
\cite[4]{rauscher_systematic_2025}).

Darüber hinaus können Caches als „write-through“ oder
„write-back“ Caches implementiert werden. Write-through
beschreibt das Korrigieren eines geschriebenen Wertes im
Speicher, sobald dieser im Cache geändert wird. Bei write-back
dagegen, wird erst bei der Eviction des Blocks der Wert in
den restlichen Speicherschichten neu gesetzt.
(\cite[526,527]{sarangi_basic_2021})
Bei mehreren Kernen kann ein Cache „private“, dass bedeutet
nur für einen Teil der Kerne zugreifbar, oder „shared“, das heißt für
alle Kerne zugreifbar, sein. Wenn ein Cache-Layer private ist,
bedeutet das außerdem, dass für jeden Kern dieser Cache
separat physisch implementiert ist.
(\cite[593,594]{sarangi_basic_2021})
Auch wenn die im folgenden
betrachteten \acp{CPU} Instruktionen und Daten in demselben
Speicher halten, ist es üblich, den L1 Cache in einen Cache für
Instruktionen (L1i) und einen für Daten (L1d) zu teilen.
(\cite[3]{schwarz_zombieload_2019}) Diese Unterteilung ist
für \ac{CbCC} weniger relevant, weswegen im folgenden alle
Caches so betrachtet werden, als würden sie sowohl Daten als
auch Instruktionen speichern.
Je nachdem, welche Architektur im Cache-Design eines Prozessors
implementiert ist, sind verschiedene \acp{CbCC} nutzbar.
(\cite[4,5]{rauscher_systematic_2025})
Eine in der Vergangenheit häufige Cache-Architektur, beinhalteten
zwei bis vier Layer mit ein bis drei private Layern und einem
shared \ac{LLC}. In Abbildung \ref{figure:cache-architectures}
sind die Cachearchitekturen der getesteten \acp{CPU} dargestellt;
darunter befindet sich auch diese Architektur.
Die verschiedenen Cache-Implementierungen sind in Anhang
\ref{appendix:multilayer-cache} detaillierter erläutert.

\begin{figure}
    \centering
    \tikzset{shape cache/.style = {
        shape=rectangle, color=black, draw,
        outer xsep=0.1cm, outer ysep=0.1cm,
    }}
    \begin{subfigure}{0.3\textwidth}
        \centering
        \begin{tikzpicture}
            \begin{scope}[on background layer={color=orange!30}, outer sep=-0.1cm]
                \fill (-1,-1) rectangle (1,1);
                \node[] (ccenter) [centered=of bg.center] {};
                \node[] (topcenter) [above=0.6cm of ccenter.north] {};
                \node[] (botcenter) [below=0.2cm of ccenter.south] {};
                \node[shape cache, fill=green!40] (l11) [left=0cm of topcenter.center] {L1};
                \node[shape cache, fill=green!40] (l12) [right=0cm of topcenter.center] {L1};
                \node[shape cache, fill=blue!40] (l21) [left=0cm of ccenter.center] {L2};
                \node[shape cache, fill=blue!40] (l22) [right=0cm of ccenter.center] {L2};
                \node[shape cache, fill=red!40] (l31) [below=0cm of botcenter.center] {LLC};
            \end{scope}
        \end{tikzpicture}
        \subcaption{Intel\textregistered\space Core\texttrademark\space i7-6700K}
    \end{subfigure}
    \begin{subfigure}{0.3\textwidth}
        \centering
        \begin{tikzpicture}
            \begin{scope}[on background layer={color=orange!30}, outer sep=-0.1cm]
                \fill (-1,-1) rectangle (1,1);
                \node[] (ccenter) [centered=of bg.center] {};
                \node[] (topcenter) [above=0.6cm of ccenter.north] {};
                \node[] (botcenter) [below=0.2cm of ccenter.south] {};
                \node[shape cache, fill=green!40] (l11) [left=0cm of topcenter.center] {L1};
                \node[shape cache, fill=green!40] (l12) [right=0cm of topcenter.center] {L1};
                \node[shape cache, fill=blue!40] (l21) [left=0cm of ccenter.center] {L2};
                \node[shape cache, fill=blue!40] (l22) [right=0cm of ccenter.center] {L2};
                \node[shape cache, fill=red!40] (l31) [below=0cm of botcenter.center] {LLC};
            \end{scope}
        \end{tikzpicture}
        \subcaption{Intel\textregistered\space Core\texttrademark\space i7-7700HQ}
    \end{subfigure}
    \begin{subfigure}{0.3\textwidth}
        \centering
        \begin{tikzpicture}
            \begin{scope}[on background layer={color=orange!30}, outer sep=-0.1cm]
                \fill (-1,-1) rectangle (1,1);
                \node[] (ccenter) [centered=of bg.center] {};
                \node[] (topcenter) [above=0.3cm of ccenter.north] {};
                \node[] (botcenter) [below=0.25cm of ccenter.south] {};
                \node[shape cache, fill=green!40] (l11) [left=0cm of topcenter.center] {L1};
                \node[shape cache, fill=green!40] (l12) [right=0cm of topcenter.center] {L1};
                \node[shape cache, fill=blue!40] (l21) [left=0cm of botcenter.center] {L2};
                \node[shape cache, fill=blue!40] (l22) [right=0cm of botcenter.center] {L2};
            \end{scope}
        \end{tikzpicture}
        \subcaption{AMD A10-7850K APU\label{figure:cache-architecure-amd-a10}}
    \end{subfigure}
    \caption{Cache-Architekturen der getesteten \acp{CPU} (nach Anhang
    \ref{appendix:lscpu-cache})\label{figure:cache-architectures}.
    Die Caches sind jeweils nur für zwei Kerne dargestellt, um das
    Konzept der Architektur zu veranschaulichen.}
\end{figure}

\subsubsection{Pipelined Prozessoren}
Um Instruktionen auszuführen werden diese in \acp{mop}
übersetzt, die sequentiell ausgeführt werden müssten.
Alle \ac{mop} durchlaufen dabei verschiedene
Phasen. Beispielsweise müssen die \acp{mop} und ihre Operanden aus
dem Speicher in die Register geladen, in der ALU
ausgeführt und das Ergebnis anschließend in ein Zielregister geschrieben
werden. Das Ergebnis dieser Architektur
sind einzelne Einheiten, die miteinander Verbunden sind und in
einer festen Reihenfolge arbeiten. In jedem Taktzyklus
der \ac{CPU} wird eine Phase einer \ac{mop}
abgearbeitet. Da die \acp{mop} allerdings sequenziell die
Phasen durchlaufen bleibt eine Menge Zeit in einzelnen
Einheiten ungenutzt. Aus diesem Grund werden Prozessoren als
Pipeline gebaut, in der mehrere, aufeinanderfolgende \acp{mop},
parallel in verschiedenen Phasen ausgeführt werden.
(\cite[361,362]{sarangi_basic_2021})

Die Idee des Pipelining besteht darin, die Einheiten der
\ac{CPU} in sequenzieller Form anzuordnen, sodass in jedem
Zyklus jede Phase eine \ac{mop} verarbeitet. Jeder Übergang
zwischen zwei Einheiten wird dabei durch Register realisiert,
die die maximal nötigen Daten für die folgenden Phasen
zwischenspeichern. Die Phasen werden bei modernen \acp{CPU} immer
granularer unterteilt, um mehr \acp{mop} gleichzeitig in kürzer
werdenden Zyklen verarbeiten zu können.
(\cite[414-416]{sarangi_basic_2021})

Ein Problem dieser Struktur besteht darin, dass \acp{mop}
Abhängigkeiten zu früheren \acp{mop} haben können und die
Verarbeitung der vorherigen \ac{mop} zumindest teilweise
abgeschlossen sein muss. Ein typisches Beispiel dafür
sind Branch-Instruktionen. Diese haben die Aufgabe den Control
Flow, also welche Programmteile ausgeführt werden, dynamisch zu
ändern. Dies geschieht, indem der \ac{PC} auf eine neue
Zieladresse gesetzt wird. Der neue Wert des \ac{PC} verweist entweder
auf einen anderen Programmteil als den aktuell ausgeführten, dann
wird der Branch „genommen“, oder der Programmablauf wird nicht geändert
und der \ac{PC} wird auf die nächste Instruktion gesetzt, in diesem
Fall wird der Branch nicht „genommen“.
(\cite[31,34]{kaeli_speculative_2005})

Die Veränderung des \ac{PC} erfolgt an
einem Punkt in der Pipeline, wo die nächsten \acp{mop}
bereits in die Pipeline geladen und verarbeiteten werden sollten.
Hinzu kommt, dass ein Prefetching-Mechanismus Instruktionen aus
dem Speicher lädt, bevor diese ausgeführt werden.
Dies ist notwendig, da das Laden der Instruktionen zum Zeitpunkt
ihrer Ausführung von der Speicherlatenz abhängen würde und
die Pipeline erhebliche Leerlaufzeiten hätte.
(\cite[9-11]{kaeli_speculative_2005})
Diese Abhängigkeit wird als Control-Hazard bezeichnet und muss
in dem Aufbau der \ac{CPU} berücksichtigt werden.

Ein sogenanntes Data-Hazard kann dagegen auftreten, wenn eine
spätere Instruktion Daten einer früheren Instruktion verwendet,
obwohl diese frühere Instruktion den Schritt des Schreibens der
Daten in ein Register noch nicht erreicht hat. Diese Art von
Abhängigkeit wird auch als \ac{RAW}-Abhängigkeit bezeichnet.
(\cite[423-427]{sarangi_basic_2021})

Bei einer Branch Instruktion müsste also der
Prefetching-Mechanismus und die gesamte Pipeline warten, bis die
Branch-Entscheidung getroffen und der \ac{PC} angepasst wurde.
Je länger die Pipeline ist und je später die
Branch-Entscheidung erfolgt, desto länger ist diese Wartezeit.
(\cite[29-36]{kaeli_speculative_2005})

\subsubsection{Transient Execution}
Die Pipeline-Struktur führt zum Ausführen
von \acp{mop}, ohne dass vorherige \acp{mop} beendet oder
alle Abhängigkeiten aufgelöst sind. Die \ac{CPU} speichert die
Zustände der einzelnen \acp{mop} und übernimmt (engl. committed)
die architektonischen Änderungen in der korrekten Reihenfolge,
sobald alle Abhängigkeiten aufgelöst wurden. Führt eine
Abhängigkeit zu einem Fehler, werden die Änderungen verworfen
(engl. squashed) und die \acp{mop} wurden nur „flüchtig“,
also transient, ausgeführt. Änderungen in dem Bereich der
Mikroarchitektur bleiben allerdings teilweise erhalten.
(\cite[251]{canella_systematic_2019})
Für das Speichern des Zustands der \acp{mop} und das Committen
nach dem Auflösen der Abhängigkeiten ist der sogenannte \ac{ROB}
zuständig. (\cite[493,494]{sarangi_basic_2021})
Da dieser Speicher begrenzt ist kann nur eine gewisse
Menge an \acp{mop} transient ausgeführt werden
(\cite[6]{mosier_analyzing_2025}).

Ein weiterer Optimierungsmechanismus, der eng mit Pipelining
zusammenhängt, ist \ac{OoO}-Execution. Ziel ist es,
Instruktionen ohne Data Hazards effizienter auszuführen.
Dazu werden voneinander unabhängige Instruktionen in
veränderter Reihenfolge ausgeführt, um die Auslastung der
\ac{CPU} Einheiten zu maximieren. Durch das Ausführen späterer
Instruktionen vor den zuvor kommenden Instruktionen führt
allerdings auch dieses Konzept zu Transient Execution.
Wenn beispielsweise eine im Programm frühere Instruktion
zeitlich nach folgenden Instruktionen ausgeführt wird und einen
Fehler auslöst, können diese späteren Instruktionen bereits
ausgeführt worden sein, obwohl sie im eigentlichen
Programmablauf erst nach dem Fehler folgen würden.
Die im Programm nachfolgenden Instruktionen wurden zum Zeitpunkt
des Fehlers bereits ausgeführt und trotz des zurücksetzen des
Zustands können Informationen durch Covert- beziehungsweise
Side-Channel übertragen werden.
(\cite[5]{lipp_meltdown_2020}, \cite[633]{cheng_speclfb_2024})

Spekulative Ausführung ist eine spezielle Form der transient
Execution, bei der die \ac{CPU} Control- oder Data-Hazards durch das
Vorhersagen von Abhängigkeiten auflöst.
Die Anzahl der Instruktionen beziehungsweise die Zeitspanne, über
die die \ac{CPU} Instruktionen ausführt, die später verworfen
werden, wird als die Größe oder Tiefe des spekulativen
beziehungsweise transient Fensters bezeichnet. Diese Zeit hängt
von der Größe des \ac{ROB} sowie der Dauer ab, bis die zugrunde
liegenden Abhängigkeiten aufgelöst werden können. Die Größe
dieses Fensters ist beispielsweise besonders groß, wenn eine
benötigte Adresse nicht im Cache vorhanden ist und aus dem \ac{RAM}
geladen werden muss. In diesem Falle ist das Fenster nur durch
die Größe des \ac{ROB} begrenzt. (\cite[6]{mosier_analyzing_2025})
Transient Execution wird am Beispiel von spekulativer
Branch-Ausführung in Anhang
\ref{appendix:speculative-branches} näher erläutert.
(\cite[136]{kaeli_speculative_2005})
